\chapter{Défis techniques et solutions}
\label{chap:defis}

La construction de PrognoSense a soulevé plusieurs difficultés concrètes, dont la résolution
constitue une part importante de la valeur d'ingénierie du projet.

\section{Hétérogénéité des données}
\textbf{Problème.} Les cinq datasets diffèrent par leur format (CSV, MAT, TXT), leur
fréquence d'échantillonnage (12 à 25\,kHz, ou des cycles pour CMAPSS) et leur nombre de
canaux (1 à 21). \textbf{Solution.} Un adaptateur unifié (\texttt{UnifiedSignal}) normalise
tout en amont (rééchantillonnage à 12{,}8\,kHz, fenêtrage, Z-score), permettant au reste du
pipeline d'ignorer la provenance des données.

\section{Valeurs manquantes de CMAPSS}
\textbf{Problème.} Le dataset CMAPSS comportait \textbf{119\,362 valeurs manquantes} (issues
de capteurs constants pour certains régimes), bloquant l'entraînement.
\textbf{Solution.} Détection systématique et remplacement par zéro après normalisation,
combinés à une sélection implicite via l'importance des features.

\section{Prévention de la fuite de données}
\textbf{Problème.} Un découpage aléatoire des fenêtres aurait réparti des extraits d'un même
enregistrement entre apprentissage et test, gonflant artificiellement les scores.
\textbf{Solution.} Un découpage \textbf{par fichier source} (\texttt{split\_by\_source})
garantit l'étanchéité, avec ajustement automatique de la taille de test pour couvrir toutes
les classes (cas CWRU : 20\,\% $\rightarrow$ 25\,\%). Les scores rapportés sont donc
honnêtes.

\section{Non-linéarité du dataset Mechanical Faults}
\textbf{Problème.} Sur Mechanical Faults, les modèles à base d'arbres plafonnent à 75\,\% et
le Decision Tree s'effondre à 51\,\%. \textbf{Solution.} Le benchmark systématique a révélé
que le \textbf{MLP} (89{,}6\,\%) exploite mieux les interactions entre les 132 features
multi-capteurs. Cela a justifié la stratégie \emph{un modèle actif par dataset} plutôt qu'un
modèle unique imposé.

\section{Réentraînement défaillant}
\textbf{Problème.} Le réentraînement initial échouait : les modèles étaient encodés en
entiers \texttt{[0,1,2,3]} tandis que l'interface envoyait un label \emph{texte}, et fournir
les seules données d'un nouveau défaut produisait une unique classe (impossible à entraîner).
\textbf{Solution.} Réécriture complète : combinaison des nouveaux échantillons avec les
données existantes du dataset, ré-encodage propre (\texttt{LabelEncoder} frais), garde-fou
exigeant au moins deux classes, et mise à jour cohérente modèle + encodeur + scaler. Le
réentraînement déclenche ensuite automatiquement le versioning.

\section{Concurrence d'accès à la base}
\textbf{Problème.} L'écriture simultanée par la boucle de simulation, la saisie d'événements
et le versioning provoquait des erreurs \og database is locked \fg.
\textbf{Solution.} Activation du mode \textbf{WAL} de SQLite et d'un \texttt{busy\_timeout},
autorisant lectures et écriture à coexister sans blocage.

\section{Scores d'anomalie instables sur un échantillon}
\textbf{Problème.} L'ensemble d'anomalie normalisait ses scores au sein du lot courant,
donnant un score nul pour un échantillon isolé (cas du temps réel).
\textbf{Solution.} Normalisation contre une \textbf{référence de normalité} (percentiles de
la fonction de décision sur données saines), rendant le score exploitable même pour une seule
mesure.

\section{Robustesse de l'interface}
\textbf{Problème.} Une erreur de rendu d'une page (affichage d'un objet non sérialisé)
faisait disparaître toute l'application (écran blanc). \textbf{Solution.} Correction du rendu
fautif et ajout d'un \emph{error boundary} global isolant les erreurs par page. Par ailleurs,
le délai d'attente des appels au copilot a été porté à 60\,s (les réponses d'un modèle de
langage pouvant dépasser le délai par défaut de 10\,s).

\section{Synthèse}
Ces difficultés, et surtout leur résolution méthodique, illustrent le passage d'un
\emph{prototype} fonctionnant dans des conditions idéales à une \emph{plateforme robuste}
tolérante aux cas réels (données imparfaites, concurrence, erreurs d'exécution, latences).
